% 子文件 3.导论，仅input导入，无document
\noparindent
\section{导论 马克思主义中国化时代化的历史进程和理论成果}
\vspace{1em}

\subsection{两个相结合}
把马克思主义基本原理同中国具体实际相结合、同中华优秀传统文化相结合。

\vspace{1em}
\subsection{马克思主义中国化时代化的提出及其原因和内涵}
\begin{enumerate}
    \item \keytext{提出}：1938年，毛泽东在党的六届六中全会上作的题为《论新阶段》的政治报告中最先提出了“马克思主义中国化”这个命题。这是党首次明确提出“马克思主义中国化”的重大命题。马克思主义中国化同时包含着马克思主义时代化的意蕴。

    \item \keytext{原因}：推进马克思主义中国化时代化，是马克思主义理论本身发展的内在要求；是解决中国实际问题的客观需要。

    \item \keytext{科学内涵}：毛泽东所说的“解决中国问题”和“创造些新的东西”，实际上指向马克思主义中国化时代化的科学内涵，具体包含三层意思：
    \begin{itemize}
        \item 运用马克思主义的立场、观点和方法，观察时代、把握时代、引领时代，解决中国革命、建设、改革中的实际问题。
        \item 总结和提炼中国革命、建设、改革的实践经验，并将其上升为理论，不断丰富和发展马克思主义的理论宝库，赋予马克思主义以新的时代内涵。
        \item 运用中国人民喜闻乐见的民族语言来阐述马克思主义，使其植根于中华优秀传统文化的土壤之中，具有中国特色、中国风格、中国气派。
    \end{itemize}
\end{enumerate}

\vspace{1em}
\subsection{准确把握马克思主义中国化时代化的要求}
\begin{itemize}
    \item 推进马克思主义中国化时代化是一个追求真理、揭示真理、笃行真理的过程。
    \item 准确把握马克思主义中国化时代化，要做到坚持马克思主义与发展马克思主义相统一。
    \item 坚持和发展马克思主义，必须同中国具体实际相结合、同中华优秀传统文化相结合。
\end{itemize}

\vspace{1em}
\subsection{马克思主义中国化时代化的四个历史进程}
\begin{notetable}
\textbf{\large 马克思主义中国化时代化四个历史进程}
\par\vspace{0.5em}
\begin{tabular}{|>{\centering\arraybackslash}m{1.6cm}|>{\centering\arraybackslash}m{3.6cm}|m{6.2cm}|}
\hline
时间 & 历史分期 & 历史成就 \\ \hline
1921-1949 年 & 新民主主义革命时期 & 创造了新民主主义革命的伟大成就，为实现中华民族伟大复兴创造根本社会条件 \\ \hline
1949-1978 年 & 社会主义革命和建设时期 & 创造了社会主义革命和建设的伟大成就，为实现中华民族伟大复兴奠定根本政治前提和制度基础 \\ \hline
1978-2012 年 & 改革开放和社会主义现代化建设新时期 & 创造了改革开放和社会主义现代化建设的伟大成就，为实现中华民族伟大复兴提供充满活力的体制保证与物质条件 \\ \hline
2012 年至今 & 中国特色社会主义新时代 & 创造新时代伟大成就，提供更完善制度、更坚实物质、更主动精神力量 \\ \hline
\end{tabular}
\end{notetable}

\vspace{1em}
\subsection{马克思主义中国化时代化的理论成果及其相互关系}
\begin{enumerate}
    \item \keytext{主要理论成果}：在百年奋斗历程中，党先后形成毛泽东思想、邓小平理论、“三个代表”重要思想、科学发展观、习近平新时代中国特色社会主义思想，一脉相承、与时俱进。

    \item \keytext{成果内在关系}：全部理论成果整体呈现\keytext{一脉相承又与时俱进}特征：
    \begin{itemize}
        \item 毛泽东思想蕴含立场观点方法，为后续理论筑牢根基；
        \item 中特理论体系立足新实践，丰富发展毛泽东思想；
        \item 全部成果均是马列主义在中国的运用发展，是长期指导思想、共同思想基础。
    \end{itemize}
\end{enumerate}